\section{Additional Experimental Methodology}\label{sec:methodology}
\subsection{Questions and Experimental Record}
The additional experiments complement the original fifteen-configuration
study. They examine (i) identification quality and agreement between
the two code--data paths, (ii) construction and functional outcomes
after data handling, (iii) architecture-specific wrapping outcomes,
and (iv) zero-shot versus two-shot LLM transformation. The original
corpus and its measurements remain those of
Sections~\ref{sec:results} and~\ref{sect:var_metrics}.

Each new campaign records distinct inputs, selection criteria,
architecture, configuration, repetitions, input/output hashes, and
engine version or snapshot. A subset selected for low historical
working rate is reported as a subset, not a new assessment of the
whole corpus. Model identifier, prompt examples, generation parameters,
retry policy, dependency versions, and cache status are retained where
applicable. A run is distinguished from attempts within that run and
from a reused artifact; run numbers do not imply controlled model seeds.

\subsection{Code--Data Quality and Downstream Outcomes}
Reference-based identification measures use independently annotated
data-byte offsets $G$ and predicted offsets $P$. With data as the
positive class, precision is $|P\cap G|/|P|$, recall is
$|P\cap G|/|G|$, and $F_1$ is their harmonic mean. Undefined cases and
unresolved reference bytes are identified explicitly. Interval-level
agreement complements byte scores, with the boundary or overlap rule
specified in the report. Raw candidates and accepted intervals are
assessed separately to distinguish classification from filtering.

Without independent annotations, the analysis is a case study of
agreement and downstream functionality. Agreement with the symbolic
method does not make that method ground truth, and byte-identical
outputs do not establish identical raw proposals. Analysis limits,
unresolved states, and computed versus cached runtimes accompany the
identification outcomes.

For $N$ planned cases, $B$ accepted binaries, and $W$ successful
functional checks, the additional assessment reports
\begin{equation}
\mathrm{BuildRate}=B/N,\qquad
\mathrm{WorkingRate}_{\mathrm{all}}=W/N,\qquad
\mathrm{WorkingRate}_{\mathrm{built}}=W/B.
\end{equation}
Counts accompany percentages. Unbuilt cases remain in the planned-case
denominator, and uncompleted tests remain distinct from completed
failures. The observable, harness version, architecture, and observation
window are stated. Connections, open ports, dialogs, and clean
termination provide different empirical checks rather than a general
proof of equivalence.

For LLM runs, initial acceptance and acceptance within the retry budget
are separated. If the functional verifier participates in output
selection, the pass rate on accepted outputs is conditional on that
selection. Exhausted and restarted runs, including their time and
cost, remain part of the record.

\subsection{Static Verdicts, Similarity, and Cost}
The published detection definition is retained for comparison. For a
completed file, $d_i=m_i/a_i$ is the malicious-verdict fraction, with
$m_i$ malicious verdicts and $a_i$ engines under the declared eligibility
rule. The new report also distinguishes the fraction of completed
files with any malicious verdict:
\begin{equation}
\overline d=\frac{1}{F}\sum_{i=1}^{F}d_i,\qquad
q=\frac{1}{F}\sum_{i=1}^{F}\mathbf{1}[m_i>0].
\end{equation}
Unsupported engines, errors, timeouts, scan dates, and pending analyses
are accounted for explicitly. Pending is not zero detection. Identical
hashes may reuse a scan while retaining their membership in several
experimental conditions. The aggregation states whether payloads,
runs, or distinct binaries receive equal weight.

The original size, entropy, and cosine measurements are retained in
Section~\ref{sect:var_metrics}. The additional LLM validation score uses
\texttt{SequenceMatcher} on normalized instruction signatures,
not cosine similarity~\cite{python_difflib}. Its matching-block ratio
is $2M/(u+v)$ for $M$ matched elements and sequence lengths $u,v$.
The NOP condition additionally removes recognized no-op operations;
a score of one under that representation does not imply byte identity.
The original file and the prepared intermediate representation are
distinguished as comparison references. Any new cosine measure states
its token vocabulary and whether comparisons are paired or pooled.
Wrapping measurements also identify the decoding view.

Distinct output hashes characterize exact diversity independently of
similarity. Size is reported as bytes and paired relative change.
Timing separates analysis, preparation, API calls, and local checks
where available. Costs include rejected attempts and state the tariff
snapshot and estimation rule. Repetitions of the same payload are
grouped for uncertainty analysis; small pilots retain case-level scope.

\subsection{Validity of the Extensions}\label{sec:extension_limits}
The original limitations remain relevant. Bounded exploration can
miss paths because of missing environment behavior or resource limits;
LLM and deterministic paths can share errors; and structural validation
does not establish full semantic preservation. Changes in the harness,
dependencies, input selection, or engine version can affect apparent
working-rate improvements. Historical and new scans also differ in
time and service conditions. Architecture-specific development,
especially x86 wrapping, is associated with its own frozen version
and assessment rather than inferred from x64 results.
\draftnote{Finalize input counts, repetitions, model configuration,
annotations, harness, scan eligibility, and aggregation from the completed
campaigns. Preserve the original measurements as the published comparison.}
